\chapter{}

\section{مقدمة}
تعد طريقة التحويل التفاضلي (DTM) وسيلة عددية وتحليلية لحل المعادلات التكاملية \cite{6, 7, 8, 9} والمعادلات التفاضلية العادية والجزئية وأنظمة المعادلات التفاضلية، حيث توفر هذه الطريقة الحل على شكل سلاسل متقاربة ذات مكونات سهلة الحساب. وبالنسبة للمعادلات التفاضلية، تقوم هذه الطريقة ببناء حل تحليلي في صورة كثير حدود، وهي تختلف عن طريقة سلسلة تايلور التقليدية ذات الرتب العالية التي تتطلب حسابات رمزية للمشتقات الضرورية لدوال البيانات، إذ تستغرق طريقة سلسلة تايلور وقتاً طويلاً في الحسابات عند التعامل مع الرتب الكبيرة، بينما تمثل DTM إجراءً تكرارياً للحصول على حلول تحليلية لسلاسل تايلور للمعادلات التفاضلية. ويمكن إيجاد تطبيقات مختلفة لهذه الطريقة في عدة أبحاث \cite{6, 7, 8, 9, 10, 11}، حيث استخدم 
\en{Jang et al}. \cite{9} تطبيق مفهوم طريقة التحويل التفاضلي ذات أحجام الشبكة الثابتة لتقريب حلول مسائل القيم الأولية الخطية وغير الخطية، كما طبق \linebreak \en{Gökdoğan et al} \cite{8}. و \en{Jang et al} \cite{9}. تقنية التحويل التفاضلي ذات أحجام الشبكة الثابتة لحل مسائل القيم الأولية من الرتب العليا. وقد استُخدمت DTM أيضاً لحل نموذج Michaelis–Menten للتفاعلات الكيميائية الحيوية، حيث يمكن استخدام طريقة التحويل لتقييم الحل التقريبي بواسطة سلسلة تايلور المحدودة ومن خلال إجراء تكراري موصوف بالمعادلات المحولة التي يتم الحصول عليها من المعادلة الأصلية باستخدام عمليات التحويل التفاضلي. وتتمثل الميزة الرئيسية لهذه الطريقة في إمكانية تطبيقها مباشرة على المعادلات التفاضلية العادية الخطية وغير الخطية دون الحاجة إلى الخطية أو التقطيع أو الاضطراب، كما تبرز ميزة هامة أخرى وهي قدرتها على تقليل حجم العمل الحسابي بشكل كبير مع الاستمرار في تقديم حل السلسلة بدقة وبمعدل تقارب سريع.

\section{وصف طريقة \en{DTM}}
من  المعروف انه اذا كانت لدينا دالة مثل $u$ قابلة للاشتقاق بإستمرارية الى ما لا نهاية، فأنه يمكن كتابة الدالة $u$ على شكل متسلسلة تايلر كالتالي
\begin{equation}
	u(x) = \sum_{k=0}^{\infty} \frac{1}{k!} \frac{d^k u(x_0)}{dx^k} (x - x_0)^k
\end{equation}
نعرف التحويل التفاضلي (Differential transform) للدالة $u$ من الرتبة $k$ ، والذي يرمز له بالرمز $U(k)$ ، كالتالي
\begin{equation}
	U(k) = \frac{1}{k!} \frac{d^k u(x_0)}{dx^k}
\end{equation}
بينما معكوس التحويل التفاضلي يكون
\begin{equation}
	u(x) = \sum_{k=0}^{\infty} U(k) (x - x_0)^k
\end{equation}

\begin{table}[h!]
	\centering
	\renewcommand{\arraystretch}{2.5}
	\begin{tabular}{|c|c|}
		\hline
		\textbf{الدالة} & \textbf{الدالة المحولة}\\
		\hline
		$z(x) = u(x) \pm v(x)$ & $Z(k) = U(k) \pm V(k)$\\
		\hline
		$z(x) = \alpha u(x)$ & $Z(k) = \alpha U(k)$\\
		\hline
		$z(x) = du(x)/dx$ & $Z(k) = (k+1)U(k+1)$\\
		\hline
		$z(x) = d^2 u(x)/dx^2$ & $Z(k) = (k+1)(k+2)U(k+2)$\\
		\hline
		$z(x) = d^m u(x)/dx^m$ & $Z(k) = (k+m)!/k! \, U(k+m)$\\
		\hline
		$z(x) = u(x)v(x)$ & $Z(k) = \displaystyle \sum_{r=0}^{k} U(r)V(k-r)$\\
		\hline
		$z(x) = \alpha x^m$ & $Z(k) = 
		\begin{cases}
			\alpha, & k = m \\
			0, & k \neq m
		\end{cases}$\\
		\hline
	\end{tabular}
	\caption{التحويل التفاضلي لبعض الدوال}
\end{table}
\noindent في التطبيقات الواقعية ، الدالة $u(x)$ تكون على شكل متسلسلة نهائية في المعادلة (3) وسوف يرمز لها بالرمز 
$u_K^{DTM}$ كالتالي
\begin{equation}
	u_K^{DTM}(t) = \sum_{k=0}^{K} U(k) (t - t_0)^k
\end{equation}


\newpage
\section*{الحل بطريقة التحويل التفاضلي}
\begin{example}
	\addcontentsline{toc}{section}{حل نموذج النمو السكاني}
	المثال الاول هو محاولة لوصف النمو السكاني هي النموذج الذي قدمه Thomas Malthus عام 1798 \cite{malthus1798}، والذي استند فيه إلى فرضية مفادها أن معدل نمو السكان يتناسب طردياً مع إجمالي عدد السكان $p(t)$ عند الزمن $t$، أي أن:
	\begin{equation}
			\left\{
		\begin{array}{l}
			\dfrac{dp}{dt} = kp\\
			p(0) = p_0
		\end{array}
		\right.
	\end{equation}
	حيث $k$ هو عامل التناسب. الحل التحليلي لهذه المعادلة التفاضلية كما في الفصل الأول مثال
	رقم \ref{ex:popexample}
	\begin{equation}
		p(t) = p_0 e^{kt}
	\end{equation}
%\subsection*{الحل بطريقة DTM}
نوجد حل المعادلة (5) مع الشرط الابتدائي
$p(0)=2$ و $k=1.1$ على الفترة الزمنية $[0,5]$ بإستعمال طريقة التحويل التفاضلي ومقارنته بالحل التحليلي
$p(t) = 2 e^{1.1t}$
	\begin{equation}
		\left\{
		\begin{array}{l}
			\dfrac{dp}{dt} = 1.1p\\
			p(0) = 2
		\end{array}
		\right.
	\end{equation}
	بتطبيق تحويل DTM نحصل على:
	\[
	(k+1)P(k+1) = 1.1\,P(k)
	\]
	ومنها:
	\[
	P(k+1) = \frac{1.1}{k+1} P(k)
	\]
	\textbf{من الشرط الابتدائي:}
	\[
	P(0)=2
	\]
	\textbf{إيجاد التكرارات:}
	\[
	P(1) = \frac{1.1}{1}P(0) = 1.1 \times 2 = 2.2
	\]
	\[
	P(2) = \frac{1.1}{2}P(1) = \frac{1.1}{2} \times 2.2 = 1.21
	\]	
	\[
	P(3) = \frac{1.1}{3}P(2) = \frac{1.1}{3} \times 1.21 = 0.4437
	\]
	\[
	P(4) = \frac{1.1}{4}P(3) = \frac{1.1}{4} \times 0.4437 = 0.1220175
	\]
	\[
	P(5) = \frac{1.1}{5}P(4) = \frac{1.1}{5} \times 0.1220175 = 0.02684385
	\]
	\textbf{الحل التقريبي حتى 5 حدود:}
	\[
	p(x) \approx \sum_{k=0}^{5} P(k)x^k
	\]
	\[
	p(x) \approx 2 + 2.2x + 1.21x^2 + 0.4437x^3 + 0.1220175x^4 + 0.02684385x^5
	\]
هذا الحل يمكن مقارنته مع مفكوك ماكلورين 
$p_K^{ML}$ للحل الدقيق $p(t) = 2 e^{1.1 t}$
\[
p_K^{ML}(t) = \sum_{k=0}^{K} \frac{d^k p}{dt^k}(0) \frac{t^k}{k!} \quad (\text{نأخذ $K=6$}) 
\]
\begin{align*}
	&\frac{d^0 p}{dt^0}(0) = p(0) = 2\\[5pt]
	&\frac{d p}{dt}(t) = 2(1.1)e^{1.1t} \;\Rightarrow\; \frac{d p}{dt}(0)=2(1.1)=2.2\\[5pt]
	&\frac{d^2 p}{dt^2}(t) = 2(1.1)^2 e^{1.1t} \;\Rightarrow\; \frac{d^2 p}{dt^2}(0)=2(1.1)^2 = 2.42\\[5pt]
	&\frac{d^3 p}{dt^3}(t) = 2(1.1)^3 e^{1.1t} \;\Rightarrow\; \frac{d^3 p}{dt^3}(0)=2(1.1)^3 = 2.662\\[5pt]
	&\frac{d^4 p}{dt^4}(t) = 2(1.1)^4 e^{1.1t} \;\Rightarrow\; \frac{d^4 p}{dt^4}(0)=2(1.1)^4 = 2.9282\\[5pt]
	&\frac{d^5 p}{dt^5}(t) = 2(1.1)^5 e^{1.1t} \;\Rightarrow\; \frac{d^5 p}{dt^5}(0)=2(1.1)^5 = 3.221\\[5pt]
	&\frac{d^6 p}{dt^6}(t) = 2(1.1)^6 e^{1.1t} \;\Rightarrow\; \frac{d^6 p}{dt^6}(0)=2(1.1)^6 = 3.5431
\end{align*}

وبالتالي:
\[
p_K^{ML}(t) \approx \sum_{k=0}^{6} \frac{d^k p}{dt^k}(0)\frac{1}{k!} t^k 
\]
\[
p_6^{ML}(t) = 2 + 2.2 t + \frac{2.42}{2!} t^2 + \frac{2.662}{3!} t^3 + \frac{2.9282}{4!} t^4 + \frac{3.221}{5!} t^5 + \frac{3.5431}{6!} t^6
\]
\[
= 2 + 2.2 t + 1.21 t^2 + 0.44367t^3 + 0.12201 t^4 + 0.026842 t^5 + (4.921 \times 10^{-3} ) t^6
 \]

\section*{الحل بطريقة RK4}
	نريد حل المعادلة التالية بطريقة RK4
	\[
	p' = 1.1\,p, \quad p(0)=2
	\]
	على الفترة \([0,5]\) مع التقسيم:
	\[
	t_0=0,\; t_1=1,\; t_2=2,\; t_3=3,\; t_4=4,\; t_5=5
	\]
	وبالتالي فإن:
	\[
	h = 1
	\]
	نعرّف:
	\[
	f(t,p) = 1.1\,p
	\]
	\textbf{الخطوة الأولى: حساب \( p_1 \)}
	\[
	k_1 = f(0,2) = 2.2
	\]
	\[
	k_2 = f\left(0.5,\, 2 + 0.5 \times 2.2 \right) = f(0.5, 3.1) = 3.41
	\]
	\[
	k_3 = f\left(0.5,\, 2 + 0.5 \times 3.41 \right) = f(0.5, 3.705) = 4.0755
	\]
	\[
	k_4 = f(1,\, 2 + 4.0755) = f(1, 6.0755) = 6.68305
	\]
	\[
	p_1 = 2 + \frac{1}{6}(2.2 + 2(3.41) + 2(4.0755) + 6.68305)
	\]
	\[
	p_1 \approx 2 + \frac{1}{6}(23.85305) \approx 5.9755
	\]
	\textbf{الخطوة الثانية: حساب \( p_2 \)}
	\[
	k_1 = 1.1(5.9755) = 6.57305
	\]
	\[
	k_2 = 1.1(5.9755 + 0.5 \times 6.57305) = 10.18823
	\]
	\[
	k_3 = 1.1(5.9755 + 0.5 \times 10.18823) = 12.16723
	\]
	\[
	k_4 = 1.1(5.9755 + 12.16723) = 19.9560
	\]
	\[
	p_2 = 5.9755 + \frac{1}{6}(6.57305 + 2(10.18823) + 2(12.16723) + 19.9560)
	\]
	\[
	p_2 \approx 17.854
	\]
وبنفس الطريقة نجد بقية القيم

\begin{table}[h]
	\addcontentsline{toc}{section}{مقارنة النتائج}
	\centering
	\renewcommand{\arraystretch}{1.5}
	\begin{tabular}{|c|c|c|c|}
		\hline
		$t$ & $p_{\text{DTM}}(t)$ & $p_{\text{RK4}}(t)$ & Exact \\ 
		\hline
		0 & 2.00000000 & 2.00000000 & 2.00000000 \\ 
		1 & 6.00833205 & 5.97567500 & 6.00833205 \\ 
		2 & 18.05002700 & 17.85434585 & 18.05002700 \\ 
		3 & 54.22527782 & 53.34588407 & 54.22527784 \\ 
		4 & 162.90173570 & 159.38883290 & 162.90173730 \\ 
		5 & 489.38368110 & 476.22793210 & 489.38386460 \\ 
		\hline
	\end{tabular}
	\caption{\centering مقارنة بين طريقة التحويل التفاضلي (DTM) وطريقة رونج-كوتا من الرتبة الرابعة (RK4) مع الحل الدقيق للمعادلة التفاضلية $p'=1.1p$}
\end{table}

\begin{figure}[H]
	\centering
	\includegraphics{plot1.jpg}
	\caption{}
\end{figure}
\newpage
\begin{figure}[H]
	\raggedleft
	\includegraphics[scale=0.9]{code1.jpg}
%	\caption{}
\end{figure}
\end{example}
\newpage
\begin{example}
	\addcontentsline{toc}{section}{حل مثال آخر}
	لنأخذ معادلة تفاضلية خطية من الدرجة الأولى
	\begin{equation}
		\frac{dy}{dx} + 2y = \cos x
	\end{equation}
	مع الشرط الابتدائي
	\begin{equation}
		y(0) = 1
	\end{equation}
	\subsection*{الحل بطريقة DTM}
		بتطبيق تحويل DTM على المعادلة نحصل على العلاقة:
		\[
		(k+1)Y(k+1) + 2Y(k) = \frac{1}{k!}\cos\left(\frac{\pi k}{2}\right)
		\]
		ومنها:
		\[
		Y(k+1) = \frac{1}{k+1}\left(\frac{1}{k!}\cos\left(\frac{\pi k}{2}\right) - 2Y(k)\right)
		\]
		\textbf{من الشرط الابتدائي:}
		\[
		Y(0)=1
		\]
		\textbf{إيجاد التكرارات:}
		\[
		Y(1) = \frac{1}{1}\left(1 - 2\times 1\right) = -1
		\]
		
		\[
		Y(2) = \frac{1}{2}\left(0 - 2(-1)\right) = 1
		\]
		
		\[
		Y(3) = \frac{1}{3}\left(\frac{-1}{2} - 2(1)\right) = \frac{1}{3}\left(-\frac{1}{2} -2\right) = -\frac{5}{6}
		\]
		
		\[
		Y(4) = \frac{1}{4}\left(0 - 2\left(-\frac{5}{6}\right)\right) = \frac{5}{12}
		\]
		
		\[
		Y(5) = \frac{1}{5}\left(\frac{1}{24} - 2\left(\frac{5}{12}\right)\right) = \frac{1}{5}\left(\frac{1}{24} - \frac{10}{12}\right) = -\frac{13}{120}
		\]
		
		\textbf{الحل التقريبي حتى 5 حدود:}
		\[
		y(x) \approx \sum_{k=0}^{5} Y(k)x^k
		\]
		\[
		y(x) \approx 1 - x + x^2 - \frac{5}{6}x^3 + \frac{5}{12}x^4 - \frac{13}{120}x^5
		\]

		
		
		\subsection*{الحل بطريقة RK4}
		
		نريد حل المعادلة التالية بطريقة رانج كتا من الرتبة الرابعة:
		\[
		y' + 2y = \cos x, \quad y(0)=1
		\]
		أي:
		\[
		y' = f(x,y) = \cos x - 2y
		\]
		على الفترة \([0,2]\) مع التقسيم:
		\[
		x_0=0,\; x_1=0.25,\; x_2=0.5,\; x_3=0.75,\; x_4=1,\dots, x_8=2
		\]
		وبالتالي فإن:
		\[
		h = 0.25
		\]
		نعرّف:
		\[
		f(x,y)=\cos x - 2y
		\]
		\textbf{الخطوة الأولى: حساب \(y_1\)}
		\[
		k_1 = f(0,1)=\cos(0)-2(1)=1-2=-1
		\]
		\[
		\begin{split}
			k_2 &= f\left(0.125,\, 1 + 0.125(-1)\right) 
			= f(0.125, 0.875) \\
			&\approx \cos(0.125)-1.75 
			\approx 0.99219-1.75 
			= -0.75781
		\end{split}
		\]
		
		\[
		\begin{split}
			k_3 &= f\left(0.125,\, 1 + 0.125(-0.75781)\right) 
			= f(0.125, 0.90527) \\
			&\approx 0.99219-1.81054 
			= -0.81835
		\end{split}
		\]
		
		\[
		\begin{split}
			k_4 &= f\left(0.25,\, 1 + 0.25(-0.81835)\right) 
			= f(0.25, 0.79541) \\
			&\approx \cos(0.25)-1.59082 
			\approx 0.96891-1.59082 
			= -0.62191
		\end{split}
		\]
		
		\[
		y_1 = 1 + \frac{0.25}{6}\big(k_1 + 2k_2 + 2k_3 + k_4\big)
		\]
		
		\[
		y_1 \approx 1 + \frac{0.25}{6}(-1 -1.51562 -1.63670 -0.62191)
		\approx 0.80107375
		\]
		\textbf{الخطوة الثانية: حساب \(y_2\)}
		\[
		k_1 = f(0.25,0.8185)=\cos(0.25)-2(0.8185)\approx -0.66809
		\]
		\[
		\begin{split}
			k_2 &= f(0.375,\, 0.8185+0.125(-0.66809))
			\approx f(0.375,0.7345)\\
			&\approx \cos(0.375)-1.469
			\approx -0.5710
		\end{split}
		\]
		\[
		k_3 = f(0.375,\, 0.8185+0.125(-0.5710))
		\approx f(0.375,0.7476)
		\approx -0.5841
		\]
		\[
		k_4 = f(0.5,\, 0.8185+0.25(-0.5841))
		\approx f(0.5,0.6725)
		\approx \cos(0.5)-1.345
		\approx -0.4674
		\]
		\[
		y_2 = 0.8185 + \frac{0.25}{6}(k_1+2k_2+2k_3+k_4)
		\approx 0.667764
		\]
			وبنفس الطريقة يتم حساب بقية القيم \(y_3,\dots,y_8\)

\begin{table}[h]
	\addcontentsline{toc}{section}{مقارنة النتائج}
	\centering
	\renewcommand{\arraystretch}{1.5}
	\begin{tabular}{|c|c|c|c|}
		\hline
		$x$ & $y_{\text{DTM}}(x)$ & $y_{\text{RK4}}(x)$ & Exact \\ 
		\hline
		0.0 & 1.00000000 & 1.00000000 & 1.00000000 \\ 
		0.2 & 0.80096416 & 0.80107417 & 0.80096416 \\ 
		0.5 & 0.66764581 & 0.66776472 & 0.66764580 \\ 
		0.8 & 0.56288255 & 0.56297237 & 0.56288140 \\ 
		1.0 & 0.46564264 & 0.46567262 & 0.46561629 \\ 
		1.2 & 0.36547279 & 0.36520463 & 0.36517687 \\ 
		1.5 & 0.25978973 & 0.25767558 & 0.25766612 \\ 
		1.8 & 0.15480771 & 0.14361880 & 0.14361720 \\ 
		2.0 & 0.07343915 & 0.02639276 & 0.02639013 \\ 
		\hline
	\end{tabular}
	\caption{\centering مقارنة بين طريقة التحويل التفاضلي (DTM) وطريقة رونج-كوتا من الرتبة الرابعة (RK4) مع الحل الدقيق للمعادلة التفاضلية 
		$y' + 2y = \cos x$}
\end{table}

\begin{figure}[H]
	\centering
	\includegraphics{plot2.jpg}
	\caption{}
\end{figure}
\newpage
\begin{figure}[H]
	\raggedleft
	\includegraphics[scale=0.8]{code2.jpg}
	%	\caption{}
\end{figure}

\end{example}

